\documentclass[10pt,a4paper]{amsart}
\usepackage[margin=19mm]{geometry}
\usepackage{amsmath,amssymb,mathtools}
\usepackage[scr=rsfs,cal=euler]{mathalfa}
\usepackage{xfrac}
\usepackage[unicode,hidelinks,bookmarks=false]{hyperref}
\newcommand{\ie}{i.e.\ }
\newcommand{\cf}{cf.\ }
\newcommand{\resp}{respectively\ }
\newcommand{\Subsection}{\subsection}
\providecommand{\nobreakdash}{\nobreak\hbox{-}}
\setlength{\emergencystretch}{2em}
\multlinegap0pt
\usepackage{amssymb}
\usepackage{mathtools}
\usepackage{enumerate}
\newtheorem{lemma}{Lemma}
\newcommand{\bC}{\mathbf{C}}
\newcommand{\bN}{\mathbf{N}}
\newcommand{\cH}{\mathcal{H}}
\newcommand{\ip}[2]{\langle{#1},{#2}\rangle} % inner product
\newcommand{\norm}[1]{\lVert{#1}\rVert} % norm
\title{On the directional growth of the resolvent norm}
\author{Horia Cornean, Henrik Garde, Arne Jensen}
\date{Source: arXiv:2602.18211v2 (CC BY 4.0)}
\begin{document}
\maketitle

\noindent\emph{Source and licence.} On the directional growth of the resolvent norm. Version arXiv:2602.18211v2 (CC BY 4.0), distributed by arXiv under the Creative Commons Attribution 4.0 International licence, http://creativecommons.org/licenses/by/4.0/, as recorded in the arXiv licence metadata for this exact version and shown on the abstract page https://arxiv.org/abs/2602.18211v2. Full licence notice: \texttt{licenses/arxiv-2602.18211-CC-BY-4.0.txt}.\par\smallskip

\noindent\textit{Editorial scope.} Counted unit: the article's lemma lemma-abg (source0.tex lines 113--127). The article's complete original proof of it is delivered with it (source0.tex lines 185--196, one \texttt{proof} environment of 12 lines, which the article places away from the statement). No mathematics is written, repaired or re-derived: the statement and the proof are the article's own lines, in the article's own notation, and no retained line is substituted or re-flowed.  No further source-local context is needed: every cross-reference of every retained line resolves inside the two delivered spans.  Nothing was omitted from any retained span, so the package carries no omission record.  No retained line contains a citation, so the wrapper carries no bibliography and no bibliography entry.  No retained line contains a figure environment, an image-inclusion command or drawing code, so no image is delivered and the package declares no asset dependency.  Editorial formatting only, in addition: an AMS \texttt{amsart} preamble that restates the article's own declarations and macros used by the retained lines, namely the environments \texttt{lemma} on separate counters, together with the standard packages those lines need.  The retained environments print in this document as Lemma 1 in order of appearance, whereas the article prints the same items with the numbering of its own full document.  The complete native source of the article as distributed in the arXiv e-print for this exact version is bundled unchanged as \texttt{sources/195020/source0.tex}.
\begin{lemma}\label{lemma-abg}
	Let $z\in\rho(A)$. There exists a sequence $\{\psi_n\}_{n\in\bN}$ in $\cH$, such that $\norm{\psi_n}=1$, $n\in\bN$, and 
	%
	\begin{equation}\label{R-lim}
		\lim_{n\to\infty}\norm{R_A(z)\psi_n}=\norm{R_A(z)}.
	\end{equation}
	%
	Furthermore, the following limits exist:
	%
	\begin{align}
		\alpha &=\lim_{n\to\infty} \ip{R_A(z)\psi_n}{R_A(z)^2\psi_n},\label{alpha}\\
		\beta &= \lim_{n\to\infty}\norm{R_A(z)^2\psi_n}^2,\label{beta}\\
		\gamma &= \lim_{n\to\infty}\ip{R_A(z)\psi_n}{R_A(z)^3\psi_n}.\label{gamma} 
	\end{align}
\end{lemma} 

\begin{proof}[Proof of Lemma~\ref{lemma-abg}]
	Let $z\in\rho(A)$. There exists a sequence $\{\phi_k\}_{k\in\bN}$ in $\cH$, such that $\norm{\phi_k}=1$, $k\in\bN$, and $\lim_{k\to\infty}\norm{R_A(z)\phi_k}=\norm{R_A(z)}$. 
	The three sequences 
	%
	\begin{align}
		&\{\ip{R_A(z)\phi_k}{R_A(z)^2\phi_k}\}_{k\in\bN}, \label{eq:seq1} \\
		&\{\norm{R_A(z)^2\phi_k}^2\}_{k\in\bN}, \label{eq:seq2} \\
		&\{\ip{R_A(z)\phi_k}{R_A(z)^3\phi_k}\}_{k\in\bN}, \label{eq:seq3}
	\end{align}
	%
	are bounded in $\bC$. Hence, there is a subsequence $\{\phi_{k_j}\}_{j\in\bN}$ for which \eqref{eq:seq1} converges, a subsequence $\{\phi_{{k_j}_m}\}_{m\in\bN}$ for which \eqref{eq:seq2} converges, and finally a subsequence $\{\phi_{{{k_j}_m}_n}\}_{n\in\bN}$ for which \eqref{eq:seq3} converges. The result therefore follows with $\psi_n = \phi_{{{k_j}_m}_n}$, leading to the existence of all three limits~\eqref{alpha}--\eqref{gamma}. 
\end{proof}

\end{document}
